% Môn: Vật lý
% Lớp: 11
% Chương: Chương I. Dao động
% Bài: L11C1 Bài 1. Dao động điều hòa · Xác định và tính toán các đại lượng vận tốc, gia tốc tại thời điểm xác định
% ===== Câu 70 =====
% ID: L11C1B1-252-DS
% Mức: TH
\dangbt{Xác định và tính toán các đại lượng vận tốc, gia tốc tại thời điểm xác định}
\begin{ex}
% ID: L11C1B1-252-DS
% Mức: TH
Từ $x=8\cos(3t+\pi/3)$ cm, lập $v(t),a(t)$ và tính tại $t=0$.
\choiceTF

\end{ex}

% ===== Câu 74 =====
% ID: L11C1B1-253-DS
% Mức: VD
\begin{ex}
% ID: L11C1B1-253-DS
% Mức: VD
Vật có $A=9$ cm, $\omega=4$ rad/s, tại $x=A/2$ và đi theo chiều dương. Tính $v,a$ và xét nhanh dần hay chậm dần.
\choiceTF

\end{ex}

% ===== Câu 78 =====
% ID: L11C1B1-254-DS
% Mức: VD
\begin{ex}
% ID: L11C1B1-254-DS
% Mức: VD
Phân tích dấu của $x,v,a$ và sự thay đổi tốc độ khi vật đi từ biên âm về vị trí cân bằng.
\choiceTF

\end{ex}

% ===== Câu 105 =====
% ID: L11C1B1-255-DS
% Mức: TH
\begin{ex}
% ID: L11C1B1-255-DS
% Mức: TH
Vật dao động theo $x=7\cos(5t)$ cm.
\choiceTF
{\True $v=-35\sin(5t)$ cm/s.}
{\True $a=-175\cos(5t)$ cm/s².}
{\True $v_{max}=35$ cm/s.}
{Gia tốc bằng 0 ở biên.}
\end{ex}

% ===== Câu 108 =====
% ID: L11C1B1-256-DS
% Mức: VD
\begin{ex}
% ID: L11C1B1-256-DS
% Mức: VD
Vật có $A=8$ cm, $\omega=2$ rad/s, tại $x=A/2$ và đi theo chiều dương.
\choiceTF
{\True Gia tốc âm.}
{\True Độ lớn gia tốc bằng $16\,\text{cm}$ /s².}
{\True Vận tốc dương.}
{\True Vật chậm dần.}
\end{ex}

% ===== Câu 111 =====
% ID: L11C1B1-257-DS
% Mức: VD
\begin{ex}
% ID: L11C1B1-257-DS
% Mức: VD
Vật đi từ biên âm về vị trí cân bằng.
\choiceTF
{\True Li độ âm.}
{\True Vận tốc dương.}
{\True Gia tốc dương.}
{Tốc độ giảm dần.}
\end{ex}

% ===== Câu 112 =====
% ID: L11C1B1-258-TLN
% Mức: TH
\begin{ex}
% ID: L11C1B1-258-TLN
% Mức: TH
Với $x=5\cos(4t)$ cm, tốc độ cực đại theo cm/s bằng bao nhiêu?
\shortans{20}
\end{ex}

% ===== Câu 115 =====
% ID: L11C1B1-259-TLN
% Mức: VD
\begin{ex}
% ID: L11C1B1-259-TLN
% Mức: VD
Tại $x=3$ cm, $\omega=5$ rad/s. Độ lớn gia tốc theo cm/s² bằng bao nhiêu?
\shortans{75}
\end{ex}

% ===== Câu 118 =====
% ID: L11C1B1-260-TLN
% Mức: VD
\begin{ex}
% ID: L11C1B1-260-TLN
% Mức: VD
Vật có $v_{max}=14$ cm/s và $A=7$ cm. Tần số góc theo rad/s bằng bao nhiêu?
\shortans{2}
\end{ex}

% ===== Câu 169 =====
% ID: L11C1B1-261-TN
% Mức: TH
\begin{ex}
% ID: L11C1B1-261-TN
% Mức: TH
Vật dao động theo $x=9\cos(4t)$ cm. Tốc độ cực đại bằng
\choice
{\True $36\,\text{cm}$ /s}
{$9\,\text{cm}$ /s}
{$4\,\text{cm}$ /s}
{$144\,\text{cm}$ /s}
\end{ex}

% ===== Câu 170 =====
% ID: L11C1B1-262-TN
% Mức: VD
\begin{ex}
% ID: L11C1B1-262-TN
% Mức: VD
Một vật nhỏ dao động điều hòa theo phương trình $x=Acos\left( 4\pi t \right)$ (t tính bằng s). Tính từ $t=0,$ khoảng thời gian ngắn nhất để gia tốc của vật có độ lớn bằng một nửa độ lớn gia tốc cực đại là
\choice
{\True $0,083\,\text{s}$.}
{$0,104\,\text{s}$.}
{$0,167\,\text{s}$.}
{$0,125\,\text{s}$.}
\loigiai{
Ta có chu kì $T=\dfrac{2\pi }{4\pi }=0,5s.$
Khoảng thời gian ngắn nhất để gia tốc của vật có độ lớn bằng một nửa độ lớn gia tốc cực đại là $\Delta t=\dfrac{T}{6}=0,083s.$
Tốc độ cực đại của chất điểm là ${{v}_{\max }}=\omega A=2.10=\text{20 cm/s}\text{.}$
}
\end{ex}

% ===== Câu 171 =====
% ID: L11C1B1-263-TN
% Mức: VD
\begin{ex}
% ID: L11C1B1-263-TN
% Mức: VD
Một vật dao động điều hòa với biên độ $\text{4 cm}$ và chu kì $\text{2 s}\text{.}$ Quãng đường vật đi được trong $\text{4 s}$ là
\choice
{$\text{64 cm}\text{.}$}
{$\text{16 cm}\text{.}$}
{\True $\text{32 cm}\text{.}$}
{$\text{8 cm}\text{.}$}
\loigiai{
Ta có $\text{t = 4 s = 2T }\Rightarrow \text{S = 2}\text{.4A = 32 cm}\text{.}$
}
\end{ex}

% ===== Câu 58 =====
% ID: L11C1B1-264-DS
% Mức: TH
\begin{ex}
% ID: L11C1B1-264-DS
% Mức: TH
Từ $x=9\cos(5t+\pi/3)$ cm, lập $v(t),a(t)$ và tính tại $t=0$.
\choiceTF

\end{ex}

% ===== Câu 61 =====
% ID: L11C1B1-265-DS
% Mức: VD
\begin{ex}
% ID: L11C1B1-265-DS
% Mức: VD
Vật có $A=5$ cm, $\omega=2$ rad/s, tại $x=A/2$ và đi theo chiều dương. Tính $v,a$ và xét nhanh dần hay chậm dần.
\choiceTF

\end{ex}

% ===== Câu 63 =====
% ID: L11C1B1-266-DS
% Mức: VD
\begin{ex}
% ID: L11C1B1-266-DS
% Mức: VD
Phân tích dấu của $x,v,a$ và sự thay đổi tốc độ khi vật đi từ biên âm về vị trí cân bằng.
\choiceTF

\end{ex}

% ===== Câu 77 =====
% ID: L11C1B1-267-TLN
% Mức: TH
\begin{ex}
% ID: L11C1B1-267-TLN
% Mức: TH
Với $x=6\cos(2t)$ cm, tốc độ cực đại theo cm/s bằng bao nhiêu?
\shortans{11/01/1900}
\end{ex}

% ===== Câu 78 =====
% ID: L11C1B1-268-TLN
% Mức: VD
\begin{ex}
% ID: L11C1B1-268-TLN
% Mức: VD
Tại $x=3.5$ cm, $\omega=3$ rad/s. Độ lớn gia tốc theo cm/s² bằng bao nhiêu?
\shortans{31/05/2026}
\end{ex}

% ===== Câu 81 =====
% ID: L11C1B1-269-TLN
% Mức: VD
\begin{ex}
% ID: L11C1B1-269-TLN
% Mức: VD
Vật có $v_{max}=32$ cm/s và $A=8$ cm. Tần số góc theo rad/s bằng bao nhiêu?
\shortans{03/01/1900}
\end{ex}

% ===== Câu 91 =====
% ID: L11C1B1-270-DS
% Mức: NB
\begin{ex}
% ID: L11C1B1-270-DS
% Mức: NB
Trong phương trình dao động điều hòa $x=A\cos(\omega t+\varphi)$, các đại lượng $\omega$, $\varphi$ và $(\omega t+\varphi)$ là những đại lượng trung gian giúp ta xác định:
\choiceTF
{\True Tần số và pha ban đầu}
{\True Tần số và trạng thái dao động}
{Biên độ và trạng thái dao động}
{Li độ và pha ban đầu}
\loigiai{
\begin{itemchoice}
\item (Đúng) Tần số và pha ban đầu. (Vì): $\omega$ là tần số góc $\Rightarrow$ liên quan đến tần số, $\varphi$ là pha ban đầu
\item (Đúng) Tần số và trạng thái dao động. (Vì): $\varphi$ là pha ban đầu, $(\omega t + \varphi)$ là pha tại thời điểm $t$
\item (Sai) Biên độ và trạng thái dao động. (Vì): $A$ là biên độ, $(\omega t + \varphi)$ là pha tại thời điểm $t$, giúp xác định trạng thái dao động
\item (Sai) Li độ và pha ban đầu. (Vì): $x$ là li độ, $\varphi$ là pha ban đầu
\end{itemchoice}
}
\end{ex}

% ===== Câu 92 =====
% ID: L11C1B1-271-DS
% Mức: NB
\begin{ex}
% ID: L11C1B1-271-DS
% Mức: NB
Phương trình dao động điều hoà là $x=5\cos \left(2\pi t+\dfrac{\pi}{3}\right)(\mathrm{~cm}, s)$.
\choiceTF
{\True Biên độ của dao động là $5\,\text{cm}$}
{Pha dao động tại thời điểm $t$ là $\dfrac{\pi}{3}$ rad}
{\True Tần số góc của dao động là $2\pi~ \mathrm{rad} / s$}
{Quãng đường vật đi được sau 2 dao động là $20\,\text{cm}$}
\loigiai{
\begin{itemchoice}
\item (Đúng) Biên độ của dao động là $5\,\text{cm}$. (Vì): Biên độ $A = 5$ cm $\Rightarrow$ mệnh đề 1 đúng
\item (Sai) Pha dao động tại thời điểm $t$ là $\dfrac{\pi}{3}$ rad. (Vì): Pha tại thời điểm $t$ là $\omega t + \dfrac{\pi}{3}$, không phải chỉ là $\dfrac{\pi}{3}\Rightarrow$ mệnh đề 2 sai
\item (Đúng) Tần số góc của dao động là $2\pi~ \mathrm{rad} / s$. (Vì): Tần số góc $\omega = 2\pi$ rad/s $\Rightarrow$ mệnh đề 3 đúng
\item (Sai) Quãng đường vật đi được sau 2 dao động là $20\,\text{cm}$. (Vì): Quãng đường đi được sau 2 dao động là $4A \cdot 2 = 40$ cm $\Rightarrow$ mệnh đề 4 sai
\end{itemchoice}
}
\end{ex}

% ===== Câu 93 =====
% ID: L11C1B1-272-DS
% Mức: NB
\begin{ex}
% ID: L11C1B1-272-DS
% Mức: NB
Một vật dao động điều hòa với phương trình $x=2\cos\!\left(2\pi t-\dfrac{\pi}{6}\right)\,(\mathrm{cm})$.
\choiceTF
{\True Biên độ dao động điều hòa của vật là $2\,\mathrm{cm}$}
{Chu kì dao động của phương trình là $2\,\mathrm{s}$}
{Chiều dài quỹ đạo $8\,\mathrm{cm}$}
{\True Li độ của vật ở thời điểm $t=1\,\mathrm{s}$ là $\sqrt{3}\,\mathrm{cm}$}
\loigiai{
\begin{itemchoice}
\item (Đúng) Biên độ dao động điều hòa của vật là $2\,\mathrm{cm}$. (Vì): $x=2\cos\!\left(2\pi t-\dfrac{\pi}{6}\right)\Rightarrow A=2\,\mathrm{cm}\Rightarrow$ mệnh đề 1 đúng
\item (Sai) Chu kì dao động của phương trình là $2\,\mathrm{s}$. (Vì): $\omega=2\pi\Rightarrow T=\dfrac{2\pi}{\omega}=1\,\mathrm{s}\Rightarrow$ mệnh đề 2 sai
\item (Sai) Chiều dài quỹ đạo $8\,\mathrm{cm}$. (Vì): Chiều dài quỹ đạo $2A=4\,\mathrm{cm}$ (không phải $8\,\mathrm{cm}$)  $\Rightarrow$ mệnh đề 3 sai
\item (Đúng) Li độ của vật ở thời điểm $t=1\,\mathrm{s}$ là $\sqrt{3}\,\mathrm{cm}$. (Vì): $x(1)=2\cos\!\left(2\pi\cdot1-\dfrac{\pi}{6}\right)=2\cos\!\left(\dfrac{11\pi}{6}\right)=2\cdot\dfrac{\sqrt{3}}{2}.$
\end{itemchoice}
}
\end{ex}

% ===== Câu 95 =====
% ID: L11C1B1-273-DS
% Mức: NB
\begin{ex}
% ID: L11C1B1-273-DS
% Mức: NB
Cho các phát biểu sau về dao động, phát biểu nào đúng, phát biểu nào sai?
\choiceTF
{\True Dao động tuần hoàn đơn giản nhất là dao động điều hòa}
{Dao động điều hòa là dao động trong đó li độ của vật là một hàm hay sin hay hàm tan theo thời gian}
{\True Phương trình được mô tả dưới dạng $x=A\cos (\omega t+\varphi)$ được gọi là phương trình dao động điều hòa}
{Đại lượng $(\omega t+\varphi)$ được gọi là pha ban đầu của dao động}
\loigiai{
\begin{itemchoice}
\item (Đúng) Dao động tuần hoàn đơn giản nhất là dao động điều hòa. (Vì): ao động điều hòa là trường hợp đơn giản nhất của dao động tuần hoàn
\item (Sai) Dao động điều hòa là dao động trong đó li độ của vật là một hàm hay sin hay hàm tan theo thời gian. (Vì): Dao động điều hòa là dao động có li độ biến thiên theo thời gian dưới dạng hàm sin hoặc cosin, không phải hàm tan
\item (Đúng) Phương trình được mô tả dưới dạng $x=A\cos (\omega t+\varphi)$ được gọi là phương trình dao động điều hòa. (Vì): Phương trình $x=A\cos (\omega t+\varphi)$ là dạng tổng quát của phương trình dao động điều hòa
\item (Sai) Đại lượng $(\omega t+\varphi)$ được gọi là pha ban đầu của dao động. (Vì): Đại lượng $\omega t+\varphi)$ là pha của dao động tại thời điểm $t$. Pha ban đầu của dao động là $\varphi$, là giá trị của pha tại thời điểm $t=0$
\end{itemchoice}
}
\end{ex}

% ===== Câu 97 =====
% ID: L11C1B1-274-DS
% Mức: NB
\begin{ex}
% ID: L11C1B1-274-DS
% Mức: NB
Xét dao động điều hòa có phương trình $x = 2\cos\!\left(8\pi t\right)\,(\mathrm{cm})$. 
	Chọn Đúng hoặc Sai cho các phát biểu sau:
\choiceTF
{\True Chu kỳ dao động là $T = \dfrac{1}{4}\,\mathrm{s}$}
{\True Tần số dao động là $f = 4\,\mathrm{Hz}$}
{Biên độ dao động là $A = 4\,\mathrm{cm}$}
{Vật dao động theo phương tròn đều}
\loigiai{
\begin{itemchoice}
\item (Đúng) Chu kỳ dao động là $T = \dfrac{1}{4}\,\mathrm{s}$. (Vì): $\omega = 8\pi \Rightarrow T = \dfrac{2\pi}{\omega} = \dfrac{2\pi}{8\pi} = \dfrac{1}{4}\,\mathrm{s}\Rightarrow$ Đúng
\item (Đúng) Tần số dao động là $f = 4\,\mathrm{Hz}$. (Vì): $f = \dfrac{1}{T} = 4\,\mathrm{Hz}\Rightarrow$ Đúng
\item (Sai) Biên độ dao động là $A = 4\,\mathrm{cm}$. (Vì): $A = 2\,\mathrm{cm}$ (không phải $4\,\mathrm{cm}$)  $\Rightarrow$ Sai
\item (Đúng) Vật dao động theo phương tròn đều. (Vì): ao động điều hòa là hình chiếu*
\end{itemchoice}
}
\end{ex}

% ===== Câu 98 =====
% ID: L11C1B1-275-DS
% Mức: TH
\begin{ex}
% ID: L11C1B1-275-DS
% Mức: TH
Trong phương trình dao động điều hòa $x=A\cos(\omega t+\varphi)$, các đại lượng $\omega$, $\varphi$ và $\omega t+\varphi$ là những đại lượng trung gian giúp ta xác định:
\choiceTF
{\True Tần số và pha ban đầu}
{\True Tần số và trạng thái dao động}
{\True Biên độ và trạng thái dao động}
{\True Li độ và pha ban đầu}
\loigiai{
\begin{itemchoice}
\item (Đúng) Tần số và pha ban đầu. (Vì): $\omega$ giúp xác định tần số và $\varphi$ xác định pha ban đầu
\item (Đúng) Tần số và trạng thái dao động. (Vì): $\varphi$ là pha ban đầu, $\omega t + \varphi$ là pha dao động tại thời điểm $t$, liên quan đến trạng thái dao động
\item (Đúng) Biên độ và trạng thái dao động. (Vì): $A$ là biên độ, $\omega t + \varphi$ là pha dao động tại thời điểm $t$, liên quan đến trạng thái dao động
\item (Đúng) Li độ và pha ban đầu. (Vì): $x=A\cos(\omega t+\varphi)$ là li độ, $\varphi$ là pha ban đầu
\end{itemchoice}
}
\end{ex}

% ===== Câu 99 =====
% ID: L11C1B1-276-DS
% Mức: VD
\begin{ex}
% ID: L11C1B1-276-DS
% Mức: VD
Chọn Đúng hoặc Sai cho các phát biểu sau về đại lượng trong dao động điều hòa:
\choiceTF
{\True Đơn vị của vận tốc trong dao động điều hòa là cm/s}
{\True Đơn vị của gia tốc là cm/s $^2$}
{Tần số góc $\omega$ có đơn vị là rad/m}
{Đơn vị của li độ $x$ là m/s}
\loigiai{
\begin{itemchoice}
\item (Đúng) Đơn vị của vận tốc trong dao động điều hòa là cm/s. (Vì): Vận tốc: đơn vị cm/s $\Rightarrow$ Đúng
\item (Đúng) Đơn vị của gia tốc là cm/s $^2$. (Vì): Gia tốc: đơn vị cm/$s^2$
\item (Sai) Tần số góc $\omega$ có đơn vị là rad/m. (Vì): Tần số góc $\omega$: đơn vị đúng là rad/s
\item (Sai) Đơn vị của li độ $x$ là m/s. (Vì): Li độ $x$: đơn vị là cm hoặc m, không phải m/s $\Rightarrow$ Sai
\end{itemchoice}
}
\end{ex}

% ===== Câu 125 =====
% ID: L11C1B1-277-TN
% Mức: TH
\begin{ex}
% ID: L11C1B1-277-TN
% Mức: TH
Với $x=8\cos(3t)$ cm, biểu thức vận tốc đúng là
\choice
{\True $v=-24\sin(3t)$ cm/s}
{$v=24\cos(3t)$ cm/s}
{$v=-72\cos(3t)$ cm/s}
{$v=8\sin(3t)$ cm/s}
\end{ex}

% ===== Câu 126 =====
% ID: L11C1B1-278-TN
% Mức: VD
\begin{ex}
% ID: L11C1B1-278-TN
% Mức: VD
Vật có $A=9$ cm và $\omega=4$ rad/s. Tốc độ cực đại bằng
\choice
{\True $36\,\text{cm}$ /s}
{$9\,\text{cm}$ /s}
{$4\,\text{cm}$ /s}
{$144\,\text{cm}$ /s}
\end{ex}

% ===== Câu 127 =====
% ID: L11C1B1-279-TN
% Mức: VD
\begin{ex}
% ID: L11C1B1-279-TN
% Mức: VD
Vật ở phía dương và đang chuyển động về vị trí cân bằng. Nhận xét đúng là
\choice
{\True vật nhanh dần, gia tốc hướng âm}
{vật chậm dần, gia tốc hướng dương}
{vật nhanh dần, gia tốc hướng dương}
{vật đứng yên}
\end{ex}
